\subsection{From average change to instantaneous change}

If the input changes from $a$ to $a+h$, then the output changes by
\[
f(a+h)-f(a).
\]
The quotient
\[
\frac{f(a+h)-f(a)}{h}
\]
is the average rate of change over that interval.

\begin{definitionbox}[Derivative at a point]
The derivative of $f$ at $a$ is
\[
f'(a)=\lim_{h\to0}\frac{f(a+h)-f(a)}{h},
\]
provided the limit exists.
\end{definitionbox}

\begin{examplebox}[The derivative of $x^2$]
For $f(x)=x^2$,
\[
\frac{(a+h)^2-a^2}{h}=2a+h,
\]
so
\[
f'(a)=2a.
\]
\end{examplebox}

\begin{interpretationbox}[Slope and local change]
The difference quotient is the slope of a secant line. As $h\to0$, the secant approaches the tangent, and the derivative gives its slope. In applications the same number represents an instantaneous rate, such as velocity or growth rate.
\end{interpretationbox}

\begin{examplebox}[Velocity]
If position is $s(t)=t^2$, then
\[
v(t)=s'(t)=2t.
\]
At time $t=3$, the instantaneous velocity is $6$ units of distance per unit of time.
\end{examplebox}
